\begin{abstract}
Serving models whose state spans GPU video RAM (VRAM), host dynamic RAM (DRAM), and NVMe is difficult because capacity pressure, asymmetric transfer bandwidth, queueing, imperfect prefetch, and synchronization interact under bursty load. We present RAMSES, a single-node hierarchical-memory orchestrator that combines GPU resource preallocation, NUMA-aware host allocation, and asynchronous block-aligned swapping. Its overlap-aware model separates resident demand from measured bidirectional transfer volume and uses two observable pressure ratios to select admission, prefetch, and eviction actions with hysteresis. On the reported two-A100 testbed, 16{,}000 request-level measurements across five models and four task types (five restarts each) reduce end-to-end inference latency by 34--41\% relative to the PyTorch configuration; measured across the full memory hierarchy (GPU, CPU, and DRAM), whole-node energy per token falls by 47\% and the energy--delay product by 70\%, while token throughput rises by 61\%. Model-load time and peak VRAM drop by 60\% and 15\%, respectively. On a named industrial defect-inspection task (MVTec~AD) RAMSES matches baseline accuracy with output equivalence above 0.999. A 72-hour parameterized synthetic arrival trace yields fewer allocation failures and lower latency drift. These results are empirical QoS measurements for single-node serving, not field deployment, worst-case execution-time, control-stability, or functional-safety evidence. The model, measurement schema, and anonymized trace generator provide a reproducible protocol for studying when cross-tier coordination helps and when capacity or storage bandwidth remains the limiting resource.
\end{abstract}

\begin{IEEEkeywords}
Industrial AI, hierarchical memory, GPU offloading, NUMA, inference serving, workload trace
\end{IEEEkeywords}
